\documentclass[11pt]{article}
\usepackage[margin=1in]{geometry}
\usepackage{times}
\usepackage{hyperref}
\hypersetup{colorlinks=true, linkcolor=black, urlcolor=blue, citecolor=black}
\title{Digital Root Is Not a Field Equation: Marko Rodin's Vortex-Based Mathematics and Coil}
\author{Flyxion \\ \small Independent Researcher}
\date{}
\begin{document}
\maketitle

\begin{abstract}
Marko Rodin has since the 1980s promoted "Vortex-Based Mathematics," a numerological system centered on the repeating digital root pattern of doubling sequences modulo nine, and a resulting toroidal coil winding pattern, the Rodin Coil, claimed to produce anomalous electromagnetic and, in some promotional contexts, energy-generating or propulsion-relevant effects distinct from and superior to conventional coil windings. This paper argues that the numerical pattern underlying Vortex-Based Mathematics is a well understood and mathematically unremarkable property of digital roots under doubling, following directly from elementary modular arithmetic rather than revealing any hidden physical structure, that no independent, peer-reviewed electromagnetic measurement has demonstrated a Rodin Coil producing field characteristics or efficiency gains beyond what conventional toroidal winding theory already predicts for a comparable geometry, and that the coil's promotional claims of energy generation or propulsion relevance have not been distinguished, in any independently published measurement, from the well understood conventional electromagnetic behavior of any similarly shaped toroidal winding.
\end{abstract}

\section{Introduction}

Marko Rodin developed and has promoted since the 1980s a system he calls Vortex-Based Mathematics, centered on the observation that the sequence of digital roots, the repeated summing of a number's digits until a single digit remains, of successive powers of two follows a specific repeating pattern, 1, 2, 4, 8, 7, 5, cycling indefinitely, which Rodin interprets as revealing a fundamental underlying vortex structure to numbers and, by extension, to physical reality, leading to a specific toroidal coil winding pattern, the Rodin Coil, based on mapping this numerical sequence onto the coil's winding path, claimed by Rodin and promoters to produce superior or anomalous electromagnetic properties relevant to energy generation and antigravity propulsion applications.

\section{The Numerical Pattern Is an Elementary and Unremarkable Property of Modular Arithmetic}

The specific digital root pattern Rodin identifies as significant is a direct and mathematically unremarkable consequence of doubling under modulo nine arithmetic, since digital roots are equivalent to values modulo nine, and the doubling sequence's repeating cycle length and specific digit pattern follow immediately from standard modular arithmetic applied to powers of two, a fact discoverable by elementary number theory with no connection to physical vortex structures, energy fields, or any established physical phenomenon; comparable, similarly patterned cycles exist for digital roots of other simple arithmetic sequences under modulo nine and are considered elementary recreational mathematics rather than evidence of an underlying physical principle, a status the promotional literature surrounding Vortex-Based Mathematics does not engage with or distinguish its own claimed significance from.

\section{What a Toroidal Coil's Electromagnetic Properties Actually Depend On}

Conventional electromagnetic theory provides a complete and well established account of how a coil's winding geometry, wire path, turn density, and overall shape determine its inductance, field distribution, and other electromagnetic properties, and toroidal coil windings specifically, including windings that produce complex, non-uniform turn spacing along the toroid's surface, are a well studied category within this conventional theory with no established exception for windings whose specific path happens to trace a pattern derived from digital root numerology rather than from conventional design optimization criteria; a winding pattern's electromagnetic behavior is fully determined by its actual physical geometry as described in conventional field theory, independent of what numerical or symbolic interpretation the designer assigns to the pattern used to generate that geometry.

\section{The Absence of Independent Measurement Distinguishing the Claimed Effects}

Promotional materials for the Rodin Coil describe a range of claimed superior properties, including reduced heat generation, increased efficiency, and effects relevant to energy generation and antigravity applications, but no independently conducted, peer-reviewed electromagnetic measurement study has been published demonstrating that a Rodin Coil's actual measured field characteristics, efficiency, or energy behavior depart from what conventional electromagnetic theory predicts for a toroidal winding of comparable overall geometry, turn count, and wire gauge, leaving the specific claimed advantages over conventional coil designs without independent confirmation distinguishing them from ordinary toroidal coil behavior already well characterized in conventional electrical engineering.

\section{The Entropic Gravity Framing}

Rodin Coil promotional claims regarding antigravity or propulsion relevance share the general absence-of-coupling-mechanism problem addressed throughout this series regarding magnetic and electromagnetic antigravity claims, since no established or proposed mechanism, precise enough to evaluate, connects a coil's winding geometry, however unconventional its numerical derivation, to gravitational curvature under either the standard or entropic account; the claim, as with several others in this series, borrows the general vocabulary of unconventional or sacred geometry without specifying a testable physical channel.

\section{Objections}

Supporters argue that the coil's specific winding pattern, even if its numerical derivation is mathematically unremarkable in the abstract, might still happen to produce a genuinely superior electromagnetic geometry worth evaluating independent of the numerological framing used to generate it, in the same way that a design arrived at through an unconventional or even mistaken reasoning process can still turn out to be practically useful. This is a fair general principle, and it is precisely the kind of claim independent, properly controlled electromagnetic measurement could in principle confirm or refute regardless of Rodin's own numerological interpretation, a measurement that, as noted above, has not been published in independently peer-reviewed form despite several decades of the coil's promotion and commercial availability.

\section{Conclusion}

The digital root pattern underlying Vortex-Based Mathematics is an elementary and unremarkable consequence of modular arithmetic rather than evidence of hidden physical structure, a toroidal coil's electromagnetic properties are fully determined by conventional field theory regardless of the numerological interpretation assigned to its winding pattern, and no independent, peer-reviewed measurement has confirmed the Rodin Coil's specific claimed advantages over conventional toroidal windings of comparable geometry, leaving the antigravity and energy-generation applications promoted for the coil without independent evidentiary support.

\end{document}
